\ssr{ЗАКЛЮЧЕНИЕ}

В результате выполнения работы была достигнута поставленная цель -- разработана модификация муравьиного алгоритма с динамическим подбором гетерогенных параметров на основе генетического алгоритма для решения задачи коммивояжёра.

В аналитической части была формализована задача коммивояжёра в терминах теории графов и нотации IDEF0, описан классический муравьиный алгоритм и выявлены его основные недостатки: медленная сходимость, попадание в локальные минимумы, склонность к застою, сильная зависимость от параметров и высокая вычислительная сложность. Были рассмотрены и классифицированы существующие модификации муравьиного алгоритма, проведён их сравнительный анализ по ключевым критериям, на основе которого обоснована перспективность использования эволюционного подхода для динамического подбора гетерогенных параметров.

В конструкторской части были сформулированы требования к разрабатываемому методу, представлена его формализация в нотации IDEF0 и детально описаны все этапы работы: инициализация, построение маршрутов муравьями с индивидуальными параметрами $\alpha$ и $\beta$, обновление феромонов, формирование новой популяции с использованием операторов селекции, скрещивания и мутации. Выбраны и обоснованы структуры данных для реализации метода. Предложена модульная архитектура, разделяющая компоненты на модули структур данных, ядра, стратегий и наблюдателей, что обеспечивает гибкость конфигурации и расширяемость метода.

В технологической части был обоснован выбор языка программирования Go, приведены ключевые фрагменты реализации, демонстрирующие структуру муравья с гетерогенными параметрами, вероятностное правило выбора следующей вершины, локальную оптимизацию 2-opt, основной цикл колонии с периодической эволюцией популяции. Проведено тестирование программной реализации, покрытие кода составило 42,2\%. Представлен пример работы на тестовой задаче \texttt{ulysses22.tsp} с визуализацией решения и эволюции параметров.

В исследовательской части проведены три экспериментальных исследования. В первом исследовании выполнен подбор оптимальной конфигурации эволюционных операторов: лучшей признана комбинация рулеточной селекции, BLX-кроссовера и гауссовской мутации. Во втором исследовании проведена оценка сходимости разработанного алгоритма в сравнении с классическим муравьиным алгоритмом: разработанный алгоритм без локальной оптимизации показал сопоставимое качество решений (разница 1–2\%), достигая его в 3.7–8 раз быстрее по числу итераций, добавление 2-opt локальной оптимизации улучшает качество решений на 6–8\%. В третьем исследовании проведена оценка устойчивости разработанного алгоритма к выбору параметров $\alpha$ и $\beta$: коэффициент вариации средней длины маршрута снижен с 27,96\% до 4,54\%, а размах -- со 1101,85 до 124,95 единиц.

Таким образом, поставленная цель работы была достигнута, а все задачи -- выполнены. Разработанный метод может быть рекомендован для практического использования в задачах, где критична скорость получения приемлемого решения, а также в условиях, когда оптимальные параметры задачи заранее неизвестны благодаря высокой устойчивости метода.
